Es sei $Y$ die disjunkte Vereinigung der $U_i$. Wir definieren auf $Y$ eine
\definitionsverweis {Äquivalenzrelation}{}{}
$\sim$, wobei wir Punkte

\mavergleichskette
{\vergleichskette
{ x_i
}
{ \in }{ U_i
}
{  }{
}
{    }{
}
{   }{
}
}
{}{}{}
und

\mavergleichskette
{\vergleichskette
{ x_j
}
{ \in }{ U_j
}
{  }{
}
{    }{
}
{   }{
}
}
{}{}{}
als äquivalent ansehen, wenn

\mavergleichskette
{\vergleichskette
{ x_i
}
{ \in }{ U_{ij}
}
{  }{
}
{    }{
}
{   }{
}
}
{}{}{,}

\mavergleichskette
{\vergleichskette
{ x_j
}
{ \in }{ U_{ji}
}
{  }{
}
{    }{
}
{   }{
}
}
{}{}{}
und

\mavergleichskette
{\vergleichskette
{ \varphi_{ji} (x_i)
}
{ = }{ x_j
}
{  }{
}
{    }{
}
{   }{
}
}
{}{}{}
ist. Die Eigenschaften einer Äquivalenzrelation sind dabei durch die Kozykelbedingung gesichert, siehe
[[Topologischer Raum/Verklebungsdatum/Äquivalenzrelation/Aufgabe|Aufgabe]].
Wir setzen

\mavergleichskettedisp
{\vergleichskette
{X
}
{ \defeq} { Y/ \sim
}
{ } {
}
{ } {
}
{ } {
}
}
{}{}{}
und versehen $X$ mit der
\definitionsverweis {Quotiententopologie}{}{.}
Die Verknüpfungen

\mathdisp {U_i \longrightarrow Y \longrightarrow X} {  }

sind die $\psi_i$, und $V_i$ sind die Bilder dieser Abbildungen. Daher liegen Homöomorphismen
\maabb {\psi_i} {U_i } { V_i
} {}
vor. Dabei ist zu

\mavergleichskette
{\vergleichskette
{ x
}
{ \in }{ U_i
}
{  }{
}
{    }{
}
{   }{
}
}
{}{}{}

\mavergleichskettedisp
{\vergleichskette
{ \psi_i (x)
}
{ \in} { V_j
}
{ } {
}
{ } {
}
{ } {
}
}
{}{}{}
genau dann, wenn

\mavergleichskette
{\vergleichskette
{ x
}
{ \in }{ U_{ij}
}
{  }{
}
{    }{
}
{   }{
}
}
{}{}{}
ist, da genau in diesem Fall $x$ mit
\mathl{\varphi_{ij}(x)}{} identifiziert wird. Daher ist

\mavergleichskette
{\vergleichskette
{ \psi_i(U_{ij})
}
{ = }{ V_i \cap V_j
}
{  }{
}
{    }{
}
{   }{
}
}
{}{}{.}
Die Kommutativität des Diagramms

\mathdisp {\begin{matrix} U_{ij}   & \stackrel{ \varphi_{ji} }{\longrightarrow} &  U_{ji}   & \\ & \!\!\! \!\! \psi_i \searrow & \downarrow \psi_j \!\!\! \!\! & \\ & &  V_i \cap V_j   & \!\!\!\!\! \!\!\!  \\ \end{matrix}} {  }

folgt ebenso.

 [[Kategorie:Latexseite]]
